\documentclass[language=english,cjk=true]{campusbeamer}
\title{Description alignment}
\author{Template check}
\begin{document}
\section{Labels}
\subsection{Width samples}
\begin{frame}{English labels}
\begin{description}[Host / HostName]
  \item[Host / HostName] TextA: explanation with enough words to wrap naturally onto a following line without changing the description column.
  \item[User] TextB: a short explanation.
  \item[IdentityFile] TextC: another explanation.
  \item[IdentitiesOnly] TextD: the final explanation.
\end{description}
\end{frame}
\begin{frame}{中文标签}
\begin{description}[主机别名与地址]
  \item[主机别名与地址] TextA: a description for the longest label.
  \item[用户] TextB: a short explanation.
  \item[认证文件] TextC: another explanation.
\end{description}
\end{frame}
\begin{frame}{Mixed labels at small size}
\small
\begin{description}[IdentityFile 密钥路径]
  \item[Host 主机名] TextA: explanation for a mixed label.
  \item[User 用户] TextB: a short explanation.
  \item[IdentityFile 密钥路径] TextC: another explanation.
\end{description}
\end{frame}
\begin{frame}{Columns and wrapped explanations}
\begin{columns}[T,onlytextwidth]
\begin{column}{.48\textwidth}
\small
\begin{description}[IdentitiesOnly]
  \item[User] TextA: this longer description wraps within its column while keeping the first-word alignment.
  \item[IdentitiesOnly] TextB: the next description.
\end{description}
\end{column}
\begin{column}{.48\textwidth}
\footnotesize
\begin{description}[本地材料路径]
  \item[来源] TextC: this description wraps in its own column.
  \item[本地材料路径] TextD: the next description.
\end{description}
\end{column}
\end{columns}
\end{frame}
\end{document}
